\documentclass[10pt,a4paper]{amsart}
\usepackage[margin=19mm]{geometry}
\usepackage{amsmath,amssymb,mathtools}
\usepackage[scr=rsfs,cal=euler]{mathalfa}
\usepackage{xfrac}
\usepackage[unicode,hidelinks,bookmarks=false]{hyperref}
\newcommand{\ie}{i.e.\ }
\newcommand{\cf}{cf.\ }
\newcommand{\resp}{respectively\ }
\newcommand{\Subsection}{\subsection}
\providecommand{\nobreakdash}{\nobreak\hbox{-}}
\setlength{\emergencystretch}{2em}
\multlinegap0pt
\usepackage{bm}
\usepackage{bbm}
\usepackage{dsfont}
\usepackage{xspace}
\usepackage{cancel}
\usepackage{mathtools}
\usepackage{xcolor}
\usepackage{amsthm}
\makeatletter
\@ifundefined{theorem}{\newtheorem{theorem}{Theorem}}{}
\makeatother
\DeclareMathOperator{\lcm}{lcm}
\newcommand{\Fq}{\mathbb{F}_{q}}
\title[Bounds and Equivalence of Skew Polycyclic Codes over Finite ]{Bounds and Equivalence of Skew Polycyclic Codes over Finite Fields}
\author{Hassan Ou-azzou, Anna-Lena Horlemann, Nuh Aydin}
\date{Source: arXiv:2507.17571v2 (CC BY 4.0)}
\begin{document}
\maketitle

\noindent\emph{Source and licence.} Bounds and Equivalence of Skew Polycyclic Codes over Finite Fields. Version arXiv:2507.17571v2 (CC BY 4.0), distributed under the Creative Commons Attribution 4.0 International licence, as recorded in the arXiv licence metadata for this exact version and shown on the abstract page https://arxiv.org/abs/2507.17571v2. Full rights notice: licenses/arxiv-2507.17571-CC-BY-4.0.txt.\par\smallskip

\noindent\textit{Editorial scope.} Counted unit: the Theorem whose source label is \texttt{Th\_general} in the article (source0.tex lines 1864--1880, source label \texttt{Th\_general}), delivered together with the article's own complete original proof of it (source0.tex lines 1882--1972, one \texttt{proof} environment of 91 lines). No mathematics is written, repaired or re-derived: statement and proof are the article's own text line for line. Required context retained uncounted: Th.1 (lines 1337--1351). Every definition, display and label of those lines is retained unchanged. Omitted from this package and not redistributed: 1--1336, 1352--1863, 1881--1881, 1973--2659. Nothing inside a retained excerpt is omitted or altered: every retained body is delivered exactly as the article's lines are written, so no omission receipt of any kind accompanies this package. Citations: the retained excerpts carry no citation command at all, so no bibliography entry of the article is transcribed and no reference is redistributed. Reference adaptation: none was needed and none was made; every reference inside the delivered unit resolves inside it, so no unresolved reference or citation is produced. Printed slips kept literal: no printed word, symbol, hypothesis or number of the counted statement or of its proof was altered. Layout and class: the article's own class is \texttt{article} with packages including amsmath, amssymb, amsthm, amsfonts, fullpage, enumitem, biblatex, authblk, hyperref, relsize, cleveref, comment, color, xcolor; the wrapper is the lane's pinned \texttt{amsart} editorial wrapper at 10pt on A4, which restates the theorem-like environments the retained text needs and adds the article's own macro definitions that the retained text needs (\texttt{\textbackslash Fq}, \texttt{\textbackslash lcm}), with extra wrapper packages (bm, bbm, dsfont, xspace, cancel, mathtools, xcolor). The delivered numbering is the unit's own. Assets: no retained span contains a figure environment, a graphics-inclusion command, a PDF-page inclusion, a table or a TikZ picture; no image file is copied into this package, the package declares no asset dependency and no image file is redistributed with it. Licence: version arXiv:2507.17571v2 of this article is distributed under the Creative Commons Attribution 4.0 International licence, as recorded in the arXiv licence metadata for this exact version and shown on the abstract page https://arxiv.org/abs/2507.17571v2. A full rights notice is delivered at licenses/arxiv-2507.17571-CC-BY-4.0.txt. Characters: every retained excerpt is pure ASCII, so the delivered engine, XeLaTeX with its default Latin Modern text font, reports no missing character.
\begin{theorem}\label{Th.1}
 Let $0<\ell<n$ be an integer,  $(a_0, a_{\ell})$ and $(b_0, b_{\ell})$ be elements of $\mathbb{F}_q^{*} \times \mathbb{F}_q^{*}$,   $\xi$ be  a primitive element of $\Fq$   and  $ \sigma$ be an automorphism of $\Fq$. The following statements are equivalent:

\begin{enumerate}
    \item $(a_0, a_{\ell}) \sim_{(n,\ell,\sigma)} (b_0, b_{\ell}),$ i.e.,  $(a_0, a_{\ell}) $ and  $ (b_0, b_{\ell})$ are Hamming \textcolor{purple}{$(n,\ell,\sigma)$} equivalent.  

\item   There exists $\alpha\in \Fq^*$ such that $$ a_0 N_n^{\sigma}(\alpha)  = b_0   \quad  \text{and} \quad   a_{\ell} N^{\sigma}_{n-\ell}(\sigma^{\ell}(\alpha))=  b_{\ell}. $$
\item  There exists $\alpha\in \Fq^*$ such that $$ (a_0, a_{\ell}) \star
\left( N_n^{\sigma}(\alpha), N_{n-\ell}^{\sigma}(\sigma^{\ell} (\alpha)) \right)= (b_0,b_{\ell}) .$$
    
 \item   $ (a_0,a_{\ell})^{-1}\star (b_0,b_{\ell})\in H_{\ell,\sigma},$ where   $H_{\ell,\sigma}$ is the cyclic subgroup of $\Fq^*\times \Fq^*$ generated by $\left(N^{\sigma}_{n}(\xi), N_{n-\ell}^{\sigma}(\sigma^{\ell} (\xi))   \right).$
\end{enumerate}
The equivalence between (1) and (4) implies that the number of $(n,\ell,\sigma)$-equivalence classes is $$N_{(n,\ell,\sigma)}:= \dfrac{ (q-1)^2}{ \lcm\left(
\frac{q-1}{\gcd([n]_r,q-1)}, \frac{q-1}{ \gcd([n-\ell]_r,q-1)}\right)} .$$
\end{theorem}

\begin{theorem}\label{Th_general}
Let $ \vec{a}=(a_0,a_1,\ldots, a_{n-1}), \vec{b}=(b_0,b_1,\ldots, b_{n-1})\in \mathbb{F}_q^{n}$ have non-zero entries in the same $m$ positions, i.e., $a_{i_j}$ and $b_{i_j}$ are non-zero  for $0\leq i_0<\dots<i_{m-1}\leq n-1$. Moreover, let $\xi $ be a primitive element of $\Fq$. Then the following statements are equivalent:
\begin{enumerate}
\item $ \vec{a} \sim_{(n,\sigma)} \vec{b}$,  i.e., $ \vec{a} $ and  $ \vec{b}$ are  Hamming $ (n,\sigma)$-equivalent. 

\item   There exists $\alpha\in \Fq^*$ such that 
 $$  a_{i_j} N^{\sigma}_{n-i_j}(\sigma^{i_j}(\alpha)) = b_{i_j}, \   \text{for  each $ j \in \{0,1,\ldots,m-1\}$,}$$   where the  $ a_{i_j}$'s and $  b_{i_j}$'s are the non-zero components of $\vec{a}$ and $\vec{b}.$
 
 \item  There exists $\alpha\in \Fq^*$ such that 
 $$ (a_{i_0}, a_{i_1},\ldots, a_{i_{m-1}}) \left( N^{\sigma}_{n-i_0}(\sigma^{i_0}(\alpha)), N^{\sigma}_{n-i_1}(\sigma^{i_1}(\alpha)) , \ldots, N^{\sigma}_{n-i_{m-1}}(\sigma^{i_{m-1}}(\alpha))\right)=  ( b_{i_0}, b_{i_1},\ldots, b_{i_{m-1}}),$$
 where $m$ is the number of the non-zero entries  of $\vec{a}.$
\item $ (a_{i_0}, a_{i_1},\ldots, a_{i_{m-1}})^{-1}\star ( b_{i_0}, b_{i_1},\ldots, b_{i_{m-1}}) \in H,$ where $H$ is the cyclic subgroup of $ (\Fq^*)^m$ generated by 
$  \left( N^{\sigma}_{n-i_0}(\sigma^{i_0}(\xi)), N^{\sigma}_{n-i_1}(\sigma^{i_1}(\xi)) , \ldots, N^{\sigma}_{n-i_{m-1}}(\sigma^{i_{m-1}}(\xi))\right),$  and   $j=0,\ldots,m-1.$
\end{enumerate}
In particular,the number of $(n,\sigma)$-equivalence classes is
$$ N_{(n,\sigma)}=\dfrac{ (q-1)^m}{ \lcm_{ j\in \{i_0,i_1,\ldots, i_{m-1} \} }\left(\frac{q-1}{ \gcd([n-j]_r,q-1)}\right)}. $$
\end{theorem}

\begin{proof}\;
\begin{itemize}
    \item[(1) $\Rightarrow$ (2):] 
    Suppose that $ \vec{a} \sim_{(n,\sigma)} \vec{b} $, then there is $\alpha \in \mathbb{F}_q^*$ such that

$$\varphi_{\alpha} : \Fq[x;\sigma] /\langle x^n - \vec{b}(x) \rangle \to \Fq[x;\sigma] /\langle x^n - \vec{a}(x) \rangle, \quad f(x) \mapsto f(\alpha x) $$
is an $\mathbb{F}_q$-vector space isometry.  It follows that

$$ 
\varphi_{\alpha}(x^{k}) = (\alpha x)^k = N^{\sigma}_k(\alpha) x^{k} \mod (x^n - \vec{a}(x)), \quad \forall  \ k = 0, 1, \ldots, n-1.
$$

As $\varphi_{\alpha}$ is an $\mathbb{F}_q$-algebra isometry and $ \varphi(x^n - \vec{b}(x)) = 0 \mod (x^n - \vec{a}(x))$, we have
\begin{equation}
\varphi_{\alpha}(x^n) =\varphi_{\alpha} (\vec{b}(x)) = b_0 + N_1^{\sigma}(\alpha) b_1 x + \ldots + N_{n-1}^{\sigma}(\alpha)b_{n-1} x^{n-1} \mod (x^n - \vec{a}(x)).
\end{equation}

On the one hand,
\begin{equation}
\varphi(x^n) = N^{\sigma}_n(\alpha) x^n = N^{\sigma}_n(\alpha) ( a_0 + a_1 x + \ldots + a_{n-1} x^{n-1}) \mod (x^n - \vec{a}(x)).
\end{equation}

Comparing term by term, we deduce that   
$$ a_i N_{n}^{\sigma}(\alpha) = N_i^{\sigma}(\alpha)b_i, \ \text{ for any }  i \in \{0, 1, \ldots, n-1\},$$ 
and so 
$$ a_i N_{n}^{\sigma}(\alpha)  N_i^{\sigma}(\alpha^{-1}) =b_i, \ \text{ for any }  i \in \{0, 1, \ldots, n-1\}.$$ 
On the other hand, for any $  i \in \{0, 1, \ldots, n-1\} , $ we have 
$$ \begin{array}{rl}
    N_n^{\sigma}(\alpha) N_{i}^{\sigma}(\alpha^{-1})  & =  N_{i+(n-i)}^{\sigma}(\alpha) N_{i}^{\sigma}(\alpha^{-1})\\
     &=   \sigma^{i}(N_{n-i}^{\sigma}(\alpha)) N^{\sigma}_{i}(\alpha) N_{i}^{\sigma}(\alpha^{-1})  \\
     &=  \sigma^{i}(N^{\sigma}_{n-i}(\alpha))\\
     & = N^{\sigma}_{n-i}(\sigma^{i}(\alpha)).
\end{array}
$$

Hence,
$$ a_i N_{n-i}^{\sigma}( \sigma^i(\alpha))  =b_i, \ \text{ for any }  i \in \{0, 1, \ldots, n-1\}.$$ 

As  $ a_{i_j}$'s and $  b_{i_j}$'s are the non-zeros components of $\vec{a}$ and $\vec{b},$  we have
$$ a_{i_j} N_{n-{i_j}}^{\sigma}(\sigma^{i_j}(\alpha)) = b_{i_j},\  \ j=0,\ldots,m-1.$$

 \item[(2) $\Rightarrow$ (3):]  is immediate. 

 \item[(3) $\Rightarrow$ (4):]
Suppose that there is $\alpha\in \Fq^{*}$ such that 
$$  (b_{i_0}, b_{i_1},\ldots, b_{i_{m-1}} )=  \left( N_{n-i_0}^{\sigma}(\sigma^{i_0}(\alpha)), 
N^{\sigma}_{n-i_1}(\sigma^{i_1}(\alpha)), \ldots, N^{\sigma}_{n-i_{m-1}}(\sigma^{i_{m-1}}(\alpha))\right)\star ( a_{i_0},  a_{i_1}, \ldots, a_{i_{m-1}}).$$ 
For $\alpha= \xi^h,$ we obtain 
$$  ( a_{i_0}, a_{i_1},\ldots, a_{i_{m-1}})^{-1}\star (b_{i_0}, b_{i_1},\ldots, b_{i_{m-1}} ) =   \left( N_{n-i_0}^{\sigma}(\sigma^{i_0}(\xi)), 
N^{\sigma}_{n-i_1}(\sigma^{i_1}(\xi)), \ldots, N^{\sigma}_{n-i_{m-1}}(\sigma^{i_{m-1}}(\xi))\right)^h .$$
It follows that $  ( a_{i_0},  a_{i_1}, \ldots, a_{i_{m-1}})^{-1}\star ( b_{i_0},  b_{i_1}, \ldots, b_{i_{m-1}}) $ belongs to the cyclic subgroup  $H$ of $ (\Fq^{*})^m$ generated by $ \left( N_{n-i_0}^{\sigma}(\sigma^{i_0}(\xi)), 
N^{\sigma}_{n-i_1}(\sigma^{i_1}(\xi)), \ldots, N^{\sigma}_{n-i_{m-1}}(\sigma^{i_{m-1}}(\xi))\right).$
   
    \item[(4) $\Rightarrow$ (1):]  
     Suppose that $ (a^{n-{i_0}}, a^{n-{i_1}} , \ldots, a^{n-i_{m-1}})^{-1}\star (b^{n-{i_0}}, b^{n-{i_1}} , \ldots, b^{n-i_{m-1}}) $ is an element of the cyclic subgroup  $H$ of $ (\Fq^{*})^m$ generated by
$$ \left( N_{n-i_0}^{\sigma}(\sigma^{i_0}(\xi)), 
N^{\sigma}_{n-i_1}(\sigma^{i_1}(\xi)), \ldots, N^{\sigma}_{n-i_{m-1}}(\sigma^{i_{m-1}}(\xi))\right) .$$
Then there exists an integer $ h$ such that 
     $$  ( a_{i_0},  a_{i_1}, \ldots, a_{i_{m-1}})^{-1}\star ( b_{i_0},  b_{i_1}, \ldots, b_{i_{m-1}}) =  \left( N_{n-i_0}^{\sigma}(\sigma^{i_0}(\xi)), 
N^{\sigma}_{n-i_1}(\sigma^{i_1}(\xi)), \ldots, N^{\sigma}_{n-i_{m-1}}(\sigma^{i_{m-1}}(\xi))\right)^h $$
    For $ \alpha= \xi^{h}, $ we obtain  that $  a_j N_{n-j}^{\sigma}(\sigma^j(\alpha))= b_j, $ for any $ j\in \{ i_0,i_1,\ldots, i_{m-1}\}.$
As in the proof of Theorem \ref{Th.1}, we verify that  $ \varphi_{\alpha}, $ defined as  
  \begin{equation}
        \begin{array}{cccc}
        \varphi_{\alpha} : & \Fq[x;\sigma]/\langle x^n - \vec{b}(x) \rangle,  & \longrightarrow & \Fq[x;\sigma] /\langle x^n - \vec{a}(x) \rangle, \\ 
        & f(x) & \longmapsto & f(\alpha x),
        \end{array}
    \end{equation} 
    is an $\Fq$-vector space  isometry with respect to the Hamming distance.
To verify that $\varphi_{\alpha}$ is an $\mathbb{F}_q$-morphism, the key idea is to use the fact that  
$$ 
a_{i_k} N^{\sigma}_{n-i_k}(\sigma^{i_k}(\alpha)) = b_{i_k}, \quad \forall k \in \{0,1,\ldots,m-1\}, \quad \text{and} \quad N_{i+j}^\sigma(\alpha) = N_i^\sigma(\alpha)\sigma^i(N_j^\sigma(\alpha)),
$$  
to  obtain 
$$ 
N_i^\sigma(\alpha)\sigma^i(N_{n-i+j}^\sigma(\alpha)) \sigma^j(a_{k}) = N^{\sigma}_{k+j}(\alpha) \sigma^{j}(b_{k}).
$$  

From this, we can conclude that  
$$ 
\varphi_{\alpha}(f(x)g(x)) = \varphi_{\alpha}(f(x))\varphi_{\alpha}(g(x)).
$$  

We gave full details of the computation in the proof of Theorem \ref{Th.1},  so we do  not repeat the same process here.
\end{itemize}

\noindent By the equivalence between (1) and (4), we deduce that the number of $(n,\sigma)$-equivalence classes  on $(\Fq^*)^m$ corresponds to the order of the group $ (\Fq^*)^m/H,$ which equals 
$$ N_{(n,\sigma)}=\dfrac{ (q-1)^m}{ \lcm_{ j\in \{i_0,i_1,\ldots, i_{m-1} \} }\left(\frac{q-1}{ \gcd(p^{r j}[n-j]_r,q-1)}\right)}=\dfrac{ (q-1)^m}{ \lcm_{ j\in \{i_0,i_1,\ldots, i_{m-1} \} }\left(\frac{q-1}{ \gcd([n-j]_r,q-1)}\right)}. $$ 


\end{proof}

\end{document}
